```latex
\documentclass[11pt]{article}
\usepackage{amsmath,amsthm,amssymb,mathtools}
\usepackage[enumitem]{enumitem}

\newtheorem{theorem}{Theorem}
\newtheorem{lemma}[theorem]{Lemma}
\newtheorem{corollary}[theorem]{Corollary}
\theoremstyle{remark}
\newtheorem{remark}[theorem]{Remark}

\begin{document}

\title{A counterexample to the derived unit for the functor $\pi^{-1}$}
\author{}
\date{}
\maketitle

\section{Statement recalled}

Let $P$ be an (integral, saturated, sharp, commutative, unital) monoid.
Let $A_P$ denote the category whose objects are monoid homomorphisms $P\to Q$
(with $Q$ again integral, saturated, sharp), and whose morphisms $P\to Q \to Q'$
are monoid homomorphisms $Q\to Q'$ commuting with the structure maps from $P$.

Let $B_P\subset A_P$ be the full subcategory on the objects
\[
P \longrightarrow P[-x]^\sharp\qquad (x\in P),
\]
where $P\to P[-x]^\sharp$ is the initial map sending $x$ to $0$.

We are given:
\begin{itemize}
\item on $A_P^{\mathrm{op}}$ the birhizomic Grothendieck topology (as in the question),
\item on $B_P^{\mathrm{op}}$ the \emph{trivial} Grothendieck topology (all presheaves are sheaves),
\item the inclusion of categories of sites $i\colon B_P^{\mathrm{op}}\hookrightarrow A_P^{\mathrm{op}}$,
which is left exact.
\end{itemize}

By the hypothesis of the question, $i$ induces (by left Kan extension and sheafification)
a functor
\[
\pi^{-1}\colon \mathrm{Sh}(B_P)\longrightarrow \mathrm{Sh}(A_P),
\]
and $\pi_*$ denotes its right adjoint.
The question asks whether for all presheaves (equivalently, sheaves) of abelian groups
$F$ on $B_P^{\mathrm{op}}$, the unit map
\[
F \longrightarrow \mathbf{R}\pi_*\,\pi^{-1}F
\]
is a quasi-isomorphism.

\section{A formal lemma on left Kan extension}

We first isolate a purely categorical observation; it does not use any property of the
birhizomic topology.

\begin{lemma}\label{lem:initial-object}
The object $\mathrm{id}_P\colon P\to P$ is an initial object of $A_P$ and of $B_P$.
\end{lemma}

\begin{proof}
Let $a\colon P\to Q$ be an object of $A_P$. A morphism in $A_P$ from $\mathrm{id}_P\colon P\to P$
to $a\colon P\to Q$ is, by definition, a monoid map $\varphi\colon P\to Q$ such that
$\varphi\circ \mathrm{id}_P = a$, i.e.\ $\varphi=a$. Hence there is exactly one morphism
$\mathrm{id}_P\to a$, so $\mathrm{id}_P$ is initial in $A_P$.

Since $B_P$ is a full subcategory of $A_P$ containing $\mathrm{id}_P$ (it corresponds to $x=0$),
the same argument shows $\mathrm{id}_P$ is initial in $B_P$.
\end{proof}

Let $j\colon B_P\hookrightarrow A_P$ be the inclusion.

\begin{lemma}\label{lem:comma-initial}
For every object $a\in A_P$, the comma category $(j\downarrow a)$ has an initial object,
namely the object $(\mathrm{id}_P \xrightarrow{} a)$.
\end{lemma}

\begin{proof}
An object of $(j\downarrow a)$ is a pair $(b,\,b\to a)$ with $b\in B_P$ and $b\to a$ a morphism in $A_P$.
For any such $(b,\,b\to a)$, there is a unique morphism $\mathrm{id}_P\to b$ in $B_P$ by
Lemma~\ref{lem:initial-object}. We must check that this gives a morphism
\[
(\mathrm{id}_P\to a)\;\longrightarrow\;(b\to a)
\]
in the comma category, i.e.\ that the triangle
\[
\mathrm{id}_P \longrightarrow b \longrightarrow a
\]
commutes in $A_P$. But commutativity is automatic because $b\to a$ is a morphism \emph{over $P$}:
by definition of morphisms in $A_P$, the composite $P\to b \to a$ equals the structure map $P\to a$,
which is exactly the arrow $\mathrm{id}_P\to a$.

Uniqueness holds because $\mathrm{id}_P\to b$ is unique.
Thus $(\mathrm{id}_P\to a)$ is initial in $(j\downarrow a)$.
\end{proof}

\begin{corollary}\label{cor:Lan-constant}
Let $\mathcal{A}$ be a cocomplete category (e.g.\ $\mathcal{A}=\mathbf{Ab}$).
For any functor $F\colon B_P\to \mathcal{A}$, the left Kan extension
$\mathrm{Lan}_j(F)\colon A_P\to \mathcal{A}$ is (canonically) the constant functor with value $F(\mathrm{id}_P)$.
In particular, if $F(\mathrm{id}_P)=0$ in $\mathbf{Ab}$, then $\mathrm{Lan}_j(F)=0$.
\end{corollary}

\begin{proof}
By the pointwise formula for left Kan extension (see, e.g., Mac~Lane,
\emph{Categories for the Working Mathematician}, Chapter~X, \S3),
for each $a\in A_P$ one has
\[
\bigl(\mathrm{Lan}_j F\bigr)(a)\;\cong\;\operatorname*{colim}_{(b\to a)\in (j\downarrow a)} F(b).
\]
By Lemma~\ref{lem:comma-initial}, the indexing category $(j\downarrow a)$ has an initial object.
A colimit over a category with an initial object is canonically isomorphic to the value at that initial object,
so the above colimit is canonically isomorphic to $F(\mathrm{id}_P)$.
Thus $\mathrm{Lan}_j F$ is the constant functor with value $F(\mathrm{id}_P)$.

If $F(\mathrm{id}_P)=0$ in $\mathbf{Ab}$, then this constant functor is the zero functor.
\end{proof}

\section{Consequence for $\pi^{-1}$ and a counterexample}

Because the topology on $B_P^{\mathrm{op}}$ is trivial, every presheaf is a sheaf, and
\[
\mathrm{Sh}(B_P)\;\simeq\;\mathrm{PSh}(B_P^{\mathrm{op}})\;\simeq\;\mathrm{Fun}(B_P,\mathbf{Ab}).
\]
By the construction in the question, $\pi^{-1}$ is obtained by applying left Kan extension
(along $i\colon B_P^{\mathrm{op}}\hookrightarrow A_P^{\mathrm{op}}$, equivalently along $j\colon B_P\hookrightarrow A_P$
after identifying presheaves on $(\cdot)^{\mathrm{op}}$ with functors on $(\cdot)$) and then sheafifying
for the birhizomic topology on $A_P^{\mathrm{op}}$.

\begin{lemma}\label{lem:pi-inv-zero}
Let $F$ be a presheaf of abelian groups on $B_P^{\mathrm{op}}$ (equivalently $F\colon B_P\to \mathbf{Ab}$)
such that $F(\mathrm{id}_P)=0$. Then $\pi^{-1}F = 0$ in $\mathrm{Sh}(A_P)$.
\end{lemma}

\begin{proof}
By Corollary~\ref{cor:Lan-constant} applied in $\mathbf{Ab}$, the left Kan extension of $F$ from $B_P$ to $A_P$
is the zero presheaf on $A_P^{\mathrm{op}}$ (equivalently, the zero functor $A_P\to\mathbf{Ab}$).
Sheafification preserves the initial object (equivalently, sheafification of the zero presheaf is the zero sheaf).
Therefore $\pi^{-1}F=0$.
\end{proof}

\begin{theorem}\label{thm:counterexample}
Assume $P$ is not the trivial monoid (so $P$ contains some element $x\neq 0$).
Then there exists a presheaf of abelian groups $F$ on $B_P^{\mathrm{op}}$
such that the unit map
\[
F \longrightarrow \mathbf{R}\pi_*\,\pi^{-1}F
\]
is \emph{not} a quasi-isomorphism. Consequently, the answer to the question is \emph{negative in general}.
\end{theorem}

\begin{proof}
Choose $x\in P$ with $x\neq 0$, and let
\[
b_x \;:=\;\bigl(P\to P[-x]^\sharp\bigr)\in B_P.
\]
Note that $b_x\neq \mathrm{id}_P$ as an object of $B_P$ because $b_x$ sends $x$ to $0$, whereas
$\mathrm{id}_P$ sends $x$ to $x\neq 0$.

Define $F\colon B_P\to \mathbf{Ab}$ on objects by
\[
F(\mathrm{id}_P)=0,\qquad F(b_x)=\mathbb{Z},\qquad F(b)=0\ \text{for all }b\in B_P,\ b\neq \mathrm{id}_P,b_x.
\]
On morphisms $f\colon b\to b'$ in $B_P$, define $F(f)$ as follows:
\begin{itemize}
\item if $b=b'=b_x$, set $F(f)=\mathrm{id}_{\mathbb{Z}}$ (for every endomorphism $f$ of $b_x$);
\item in all other cases, $F(f)$ is the unique homomorphism between the corresponding groups, which is
necessarily the zero map because at least one of the source/target groups is $0$.
\end{itemize}
This clearly defines a functor: identity morphisms go to identity maps, and compositions are preserved
because all endomorphisms of $b_x$ are sent to $\mathrm{id}_{\mathbb{Z}}$ and all other morphisms
are sent to $0$.

By construction $F(\mathrm{id}_P)=0$ and $F\neq 0$ (since $F(b_x)=\mathbb{Z}$). Hence
Lemma~\ref{lem:pi-inv-zero} gives $\pi^{-1}F=0$ in $\mathrm{Sh}(A_P)$, and therefore
\[
\mathbf{R}\pi_*\,\pi^{-1}F \;=\; \mathbf{R}\pi_*(0)\;=\;0
\]
in the derived category $D(\mathrm{Sh}(B_P))$.

Thus the unit map $F\to \mathbf{R}\pi_*\pi^{-1}F$ is just the zero map $F\to 0$, which cannot be a quasi-isomorphism
because $F$ is a nonzero complex concentrated in degree $0$ (equivalently, $H^0(F)=F\neq 0$ whereas
$H^0(0)=0$).
\end{proof}

\begin{remark}
The counterexample above does \emph{not} use any property of the birhizomic topology; it fails already at the level
of the left Kan extension step, because $B_P$ contains the initial object $\mathrm{id}_P$ of $A_P$, making the
comma categories $(B_P\downarrow a)$ have initial objects for all $a\in A_P$.
\end{remark}

\end{document}
```